% ============================================================
%  Niels Treschow, "Om Gud, Idee- og Sandseverdenen samt de førstes
%    Aabenbarelse i den sidste. Et philosophisk Testament, tilligemed
%    Forfatterens Levnet som Fortale." Første Bog.
%  Christiania: Chr. Grøndahl, 1831.
%  TRANSCRIPTION (Danish) — SELECTION: two spans from vol. I,
%    printed pp. 202–207 and 220–221.
%
%  WHY THIS SELECTION. It settles a problem raised by Almindelig Logik § 31
%  (1813), where Treschow denies that the highest being is "noget Heelt, som
%  Verden, hvoraf alle Individuer ere Dele" and blames pantheism and the
%  emanation-system on that confusion. Read alone, that seems to cut against
%  the monism of the 1810 essay (p. 251, "enhver er dette Hele selv under sin
%  særegne Form"). These pages show the two are consistent: the unity of the
%  ideal world is constituted by VEXELVIRKNING among Ideas, not by composition,
%  and it is the SENSE-world — not God — that individuals compose.
%  Span A (pp. 202–207) states the doctrine of individual Ideas; span B
%  (pp. 220–221) states the whole-by-reciprocity.
%  Companions: ../enslige-ting/ (1810 essay), ../almindelig-logik/ (§§ 26–35).
%
%  Scan: ~/bibliotek/Treschow, Niels/om-gud-idee-og-sandseverdenen-1-1831.pdf
%        (379 scan pp.; nb.no IIIF, URN:NBN:no-nb_digibok_2008040213001;
%         public domain, "Tilgang for alle").
%  OFFSET (VERIFIED at both ends): scan page = printed + 54.
%        p. 202 = scan 256; p. 207 = scan 261; p. 220 = scan 274; p. 221 = scan 275.
%        The large offset is the prefixed autobiography ("Forfatterens Levnet
%        som Fortale"). Re-verify before working elsewhere in the volume.
%  TYPE: FRAKTUR. Emphasis = letterspacing (Sperrsatz) -> \emph{}.
%  ORTHOGRAPHY: uses "ø" (as in Almindelig Logik, not the "ö" of the 1810
%        antiqua essay). Note the later spelling: "Idee", "Sandseverden",
%        "Gjenstand"/"Gjendrivelse" with gj-, "Slægter", "Bevidsthed",
%        "Frembringelse", "Grundformen". Preserve as printed.
%  See RESUME-NOTES.md for the term map used to locate these spans.
% ============================================================
\documentclass[12pt,a4paper]{article}
\usepackage[utf8]{inputenc}
\usepackage[T1]{fontenc}
\usepackage{libertinus}
\usepackage{libertinust1math}
\usepackage[danish]{babel}
\usepackage[protrusion=true,expansion=true]{microtype}
\usepackage{setspace}
\usepackage[margin=1.2in]{geometry}
\usepackage{fancyhdr}
\usepackage{hyperref}

\hypersetup{
  pdftitle={Om Gud, Idee- og Sandseverdenen (1831) — uddrag},
  pdfauthor={Niels Treschow},
  hidelinks
}

\pagestyle{fancy}
\fancyhf{}
\rhead{Treschow, \emph{Om Gud} I (1831), uddrag}
\cfoot{\thepage}

\setstretch{1.3}

\newcommand{\spanhead}[1]{%
  \bigskip\begin{center}\rule{0.18\linewidth}{0.4pt}\\[6pt]
  {\small\itshape #1}\\[4pt]\rule{0.18\linewidth}{0.4pt}\end{center}\smallskip}

\title{\textbf{Om Gud, Idee- og Sandseverdenen}\\[4pt]
  \large samt de førstes Aabenbarelse i den sidste\\[6pt]
  \normalsize Et philosophisk Testament. Første Bog.\\[8pt]
  \large Uddrag: trykte s.\ 202--207 og 220--221}
\author{Niels Treschow}
\date{\small Christiania: Chr.\ Grøndahl, 1831}

\begin{document}
\maketitle
\thispagestyle{fancy}

\bigskip

\spanhead{Uddrag A: trykte s.\ 202--207 (scan 256--261).\\
  Ideernes individualitet; Slægts- og Artsbegrebernes ubestemthed.}

% printed p. 202
\noindent
[\dots] der uden Compas eller synlige Stjerner tumles omkring paa Havet, og
derfor gjerne lande ved den nærmeste Strandbred, hvortil Lykken fører dem.
Skulde der da ikke være noget Princip, der paa denne Undersøgelses Reise kan
bestemme, og tillige, som en fast Pol, lede vor Kurs? Er det vel usandsynligt,
og ikke snarere en Forestilling, som Sagens egen Natur maa medføre, at,
ligesaavel som det høieste Væsen nødvendig er virksomt i Frembringelsen af de
utallige Former eller Ideer, hvori det aabenbarer sig, saa maae disse af en lige
Nødvendighed gjøre det Samme? Nu kunne de, som virkelige Ting, ei være noget
andet end individuelle; thi ikkun Individuer, ikke Slægter og Arter ere i enhver
Henseende saa bestemte, at de for sig kunne være til. Skulle da Ideerne ---
hvilket man af deres Naturs Fuldkommenhed og nære Forvandtskab med den
guddommelige maa formode --- i Efterligning af denne ogsaa kunne skabe; saa maa
følgelig denne Skabning ligeledes være noget Individuelt, nemlig en Form af dem
igjen, der bestandig udvikler sig, ei flere forskjellige Former; saavidt kan
deres Kraft, som, skjøndt i en Henseende uendelig, dog selv ikkun er en enslig
Form, ikke strække sig. I et saadant Individ er Grundformen, d.~e.\ dets Væsen,
alene bestandig og uforanderlig, i de tilfældige Former, det under sin Fremgang
efterhaanden antager, kan der ingen Bestandighed være; thi disse ere kun i Tiden
mulige, hvor
% printed p. 203
idel Afvexling hersker. Saaledes opstaaer en Sandseverden af den ideale, og det
med samme Nødvendighed som den, med hvilken denne selv af en absolut Enhed og
Treenighed er oprunden.

Enhver Idee frembringer altsaa sit eget Billede, men i utallige øieblikkelig
vexlende Gestalter fra Kimen af indtil den modne Frugt. Ei engang hermed er al
videre Udvikling forbi. Døden selv, der følger paa Frugten, er selv kun en
Forvandling. Hos de perennerende Planter fremdriver ethvert Foraar friske Skud.
Træet voxer saa længe det lever, og naar det omsider døer, fortsætter det i sin
Afkom en i mange Stykker forandret, maaskee tilsidst endog forædlet Tilværelse.
Det samme er om bløde Urter og aarlige Væxter at formode. Ideen, hvoraf alt
Sandseligt er et Udtryk, er selv uendelig, men Billedet endeligt, saa at dette,
hvor meget det end haster, hvor flittigt det arbeider, dog aldrig naaer sin
Stræbens Endemaal. Vi alle denne Sandseverdens Indvaanere forholde os til vore
Ideer, ligesom disse til det absolut Ene. Begge Slags Væsener ere Udflydelser af
det Samme, men paa en forskjellig, nemlig deels umiddelbar, deels middelbar
Maade. Saaledes lægger ogsaa den menneskelige Forstand sine Planer; dens Ideer,
hvilke kunne være af utallige Slags, saasom af veldyrkede Marker, Lysthauger,
Pragtbygninger, videnskabelige Theorier, Kunstverker, offentlige og private
Indretninger til Livets Nytte, ere hver
% printed p. 204
for sig i Ophavsmandens Aand paa engang til; men udvortes blive de først meget
langsomt bragt den Fuldkommenhed nogenledes nær, som man i Begyndelsen kun
dunkelt, efterhaanden, som Værket skrider frem, stedse klarere beskuer.

Vil man sige, at det her er Forstanden, ei Ideen, som frembringer eller skaber,
saa er Forstanden, skjøndt man gemeenlig synes at ville fremstille den som en
virkelig Kraft, dog intet andet end et abstract og almindeligt Begreb, hvorunder
utallige Kraftyttringer eller egentlige Kræfter blive indbefattede. Disse ere nu
de menneskelige Ideer, og de mange Slags høiere Forestillinger, som deraf
opstaae: de samme ere følgelig ogsaa i Sandhed alene productive.

Ved alt dette maa det bemærkes, at jeg nødes til her at kalde de høiere
Forestillinger, der betragtes som Fostre af Fornuft, men, egentlig talt, udvikle
sig af en eneste, ligesom af en ædel Spire, saasom de nys anførte, med samme Navn
som hine det absolut Enes evige Former. Analogien mellem Gud og
Menneskeaanden, det Uskabte og Skabte, det Uendelige og Endelige berettiger til
denne Talebrug.

Jeg sagde nylig, at ingen evig Idee frembringer mere end sit eget individuelle
Billede. Herimod synes det at stride, at den menneskelige Forstand hos samme
Person eller Individ alligevel kan frembringe mange ganske forskjellige Planer og
Værker, der ikke kunne betragtes som blotte Former eller Billeder af ham
% printed p. 205
selv. Saaledes er et Kunststykke, ja endog den dertil i Tanken lagte Plan intet
Billede af Kunstnerens Aand, og var end en af de mange derfor at ansee, saa kunde
det dog ikke gjælde om de øvrige, der ofte have meget liden Lighed med ham eller
med hverandre. Ikke desmindre veed enhver \emph{Charactermaler}, enten han med
Penselen, i Vox, Marmor eller med Ord og Tale vil skildre nogen Person, at hans
hele Udseende, alle hans Ideer, Planer og Handlinger ere de enslige Træk, der
tilsammen udgjøre hans Individualitet; og hvis noget deraf mangler, er denne i sin
Heelhed ei saa fremstillet som den bør være. Ei heller maa det Forskjellige i
Værkernes eller de frembragte Gjenstandes Indhold herved tages i Betragtning, men
en vis umiskjendelig Lighed i Formen, hvorved alt det udmærker sig, som er samme
Individ at tilskrive. Fremdeles er Individualiteten sammensat af mange saavel
indvortes som udvortes fornemmelige Stykker. Det maa i ethvert Tilfælde være
ligesaa eller mere vanskeligt at udlede disse af nogen Grundformens Enhed, som det
er i den menneskelige Natur overhoved at udfinde det Grundvæsentlige, paa hvilket
alle ligeledes bestandige Egenskaber, som deri yttre sig, beroe; ikke desmindre er
der ingen Tvivl om, at enhver Enslig, hvor stor Mangfoldighed man hos ham
saavelsom hos hans evige Idee kan opdage, dog ikkun er en. Den Indvending, at der
fra en Idee kan udgaae utallige, der ligesaalidt
% printed p. 206
ere blotte Former af det samme Individ, som menneskelige Planer og Værker af den
individuelle Forstand, lader sig altsaa dermed besvare, at det sidste dog virkelig
er saa, og at i de forskjellige Individuer, i hvis Frembringelse denne Forstand
har nogen Deel, kun visse Træk tilhøre den, som tilsammentagne udgjøre ikke flere,
men blot et eneste Billede af den selv. Ethvert saadant Træk er en besynderlig,
men ei ved sig selv bestaaende Form; denne bestaaer alene i og ved de Gjenstande,
hvori den udtrykkes. Imellem den menneskelige Forstand og de evige Ideer er der
altsaa den betydelige Forskjel, at den første intet for sig Bestaaende kan
frembringe; hvad der giver dens Virkninger en vis Bestandighed er ikke til ved
den, men var tillige med den, nemlig det man sædvanlig kalder Tingenes Materie.
Ideerne kunne derimod, ligesom Gud, og formedelst en dem dertil af ham forundt
Virksomhed frembringe alt hvad man i Sandseverdenen anseer for selvstændigt eller
en virkelig Substants.

Den her fremsatte Theorie adskiller sig derved væsentlig fra den platoniske, at
den sidste gjør Ideerne til de almindelige Begrebers Gjenstande, eller omvendt
disse til Ideer, hvortil ei alene de Slægter og Arter maae henregnes, under hvilke
alle naturlige Ensligvæsener, saasom enslige Chrystaller, Planter og Dyr, men
endog blotte Kunstproducter, saasom Borde, Stole, o.~s.~v.\ indbefattes.
\emph{Plato} ansaae nemlig alle
% printed p. 207
Individuer for Ting, hvori der er intet ganske bestandigt, og vidste ei af nogen
uforanderlig Grundform, der efterhanden udvikler sig, hos dem at sige; endskjøndt
det er klart, \emph{Socrates} og \emph{Alcibiades} ligesaavel i al Evighed
vedblive at være Socrates og Alcibiades som Dyret et Dyr, Planten en Plante, ja,
de første --- hvilket jeg i en egen Afdeling af følgende Bog nærmere skal oplyse
--- saa meget mere, som alle \emph{Classe-, Slægts- og Arts-Begreber ere mere
eller mindre vaklende og ubestemmelige, men Individuerne i sig selv saa fuldkommen
bestemte, at en Forandring eller Overgang fra det ene til det andet er umulig},
da et Individ derimod meget vel kan gaae over til en anden Art, Slægt, Orden eller
Classe end den, hvorhen det fortiden hører. Der gives enten ingen Ideer, hvis
Tilværelse dog er beviist, eller disse bestaae i individuelle, evige,
uforanderlige og fuldkomne Grundformer, som fra den ideale Verden i
Sandseverdenen udtrykkes og stedse mere udvikles. Ere nu Ideerne individuelle, saa
kan der mellem dem ei være anden Forskjel end i Henseende til denne deres Form,
men i Fuldkommenhed maae alle, enhver efter sin Art, være lige. De staae altsaa ei
paa forskjellige Trin af Udvikling, ei heller ere de i samme Mening lavere og
høiere, som naar Slægter, Arter og Individuer siges at staae i saadant Forhold til
hverandre. Derimod maa der blandt Ideerne, ligesom i Sandseverdenen, baade gives
Indi\-[viduer \dots]

\spanhead{Uddrag B: trykte s.\ 220--221 (scan 274--275).\\
  Den ideale Verdens Enhed ved Vexelvirkning; Individerne som Sandseverdenens
  bestanddele.}

% printed p. 220
\noindent
[\dots] I den ideale maa ligesom i den reale Verden en bestandig Vexelvirkning
have Sted: den kunde ellers ingen virkelig Enhed, intet Heelt udgjøre. Dens
Medlemmer maae kunne meddele sig, kjende hverandre, uden at deres Tilstand derfor
bliver forandret; thi de leve ei i Tiden, hvori al Omskiftelse foregaaer. Denne
Meddelelse er mere eller mindre middelbar efter det naturlige Forhold, hvori de
staae indbyrdes som lige og ulige, over- og underordnede. Dette Forhold er
bestandigt; deraf reiser sig tillige en Forskjel paa Virksomhed med Hensyn til
nærmere og fjernere Gjenstande. Forsaavidt adskille disse Egenskaber hos Ideerne
sig fra den guddommelige Almagt
% printed p. 221
og Alvidenhed, i hvis Uendelighed Middelbarhed og Afstand aldeles forsvinde. Et
lige saadant Forhold vil man kunne bemærke mellem de menneskelige Fornuftideer; de
ere allesammen hverandre mere og mindre forvante; af hvilkensomhelst Art, moralske,
politiske, physiske, oeconomiske, ere de ei uden vexelviis Indflydelse, og
udgjøre, engang tilbørlig ordnede, et sammenhængende Heelt. Af os selv, af vor egen
Bevidsthed, som af den eneste os tilgængelige Kilde, maae vi hente alt, hvad vi om
den høiere Verden attraae og kunne haabe at faae at vide. Efter
Identitets-Systemets Principier ville vi ei heller deri blive bedragne; thi
væsentlig ere de høieste og laveste Ting ganske lige. Men, foruden den fælleds
Virksomhed, besidder enhver Idee tillige sin særegne, som bestaaer i Frembringelse
og Udvikling af et til dens individuelle Væsen svarende Billede under Tidens og
Rummets idelig afvexlende Former. Disse ere de Individuer, hvoraf Sandseverdenen
bestaaer. Ideerne danne sig paa denne Maade et Slags Legemer, ved hvilke, dog uden
deraf, som af noget Udvortes, at afficeres og forandres, de sætte sig istand til
med fuld Bevidsthed fra alle mulige Sider at beskue sit eget Indre, fordi de deri
kunne forfølge det gjennem alle Trin af en mulig Tilværelse, ei alene som det er
paa det allerhøieste, paa hvilket de virkelig staae. I ethvert sandseligt Individ
aabenbarer sig saaledes dets Idee som i en den eiendommelig stedse passende
[\dots]

\begin{center}\rule{0.25\linewidth}{0.4pt}\end{center}

\end{document}
